\documentclass{article}\usepackage{amsmath}\usepackage{graphicx}\usepackage{geometry}\geometry{a4paper, margin=1in}\title{SelAction Program: A Technical Description}\author{Generated by Gemini}\date{\today}\begin{document}\maketitle\begin{abstract}This document provides a technical description of the SelAction program, a tool for predicting genetic response and rates of inbreeding in animal breeding programs. The program is written in Fortran and supports a variety of selection schemes, including discrete and overlapping generations. This document details the mathematical models and equations implemented in the software.\end{abstract}\tableofcontents\section{Introduction}SelAction is a comprehensive software package designed for animal breeders to evaluate the consequences of different selection strategies. It calculates the expected genetic gain and the rate of inbreeding for complex breeding programs involving multiple traits and stages of selection. The program is divided into two main components: one for populations with discrete generations and another for populations with overlapping generations.\section{Discrete Generations}The `seldiscrete.f90` module is the core of the discrete generation models. It supports one, two, and three-stage selection.\subsection{Selection Index Theory}The program is based on the selection index theory developed by Smith (1936) and Hazel (1943). The selection index ($I$) is a linear combination of phenotypic information sources:\begin{equation}I = \mathbf{b'x}\end{equation}where $\mathbf{b}$ is a vector of index weights and $\mathbf{x}$ is a vector of phenotypic values for the various information sources. The goal is to find the vector $\mathbf{b}$ that maximizes the correlation between the index and the breeding goal, $H$. The breeding goal is a linear combination of true breeding values ($g_i$) weighted by their economic values ($v_i$):\begin{equation}H = \sum_{i=1}^{n} v_i g_i = \mathbf{v'g}\end{equation}The vector of optimum index weights is given by:\begin{equation}\mathbf{b} = \mathbf{P}^{-1}\mathbf{Gv}\end{equation}where:\begin{itemize}\item $\mathbf{P}$ is the phenotypic covariance matrix of the information sources in the index.\item $\mathbf{G}$ is the genetic covariance matrix between the information sources in the index and the traits in the breeding goal.\item $\mathbf{v}$ is the vector of economic values.\end{itemize}\subsection{Bulmer Equilibrium}The program accounts for the "Bulmer effect," which is the reduction in genetic variance due to selection. The program iteratively calculates the genetic and phenotypic (co)variance matrices until they reach an equilibrium, where the reduction in variance due to selection is balanced by the generation of new variance. The change in the additive genetic covariance between traits $i$ and $j$ is given by:\begin{equation}\Delta \sigma_{A_{ij}} = -k \frac{\text{cov}(A_i, I) \text{cov}(A_j, I)}{\sigma_I^2}\end{equation}where $k$ is the variance reduction coefficient, which is a function of the selection intensity. The program iterates until the change in the covariance matrix is negligible.\subsection{Multi-Stage Selection}For multi-stage selection, the program calculates the genetic gain from each stage of selection and sums them. The selection is sequential, and the genetic and phenotypic parameters are updated after each stage to account for the selection that has already occurred. For a two-stage selection process, the total response is:\begin{equation}\Delta G_{total} = \Delta G_1 + \Delta G_2\end{equation}where $\Delta G_1$ is the genetic gain from the first stage, and $\Delta G_2$ is the additional gain from the second stage, conditional on the selection in the first stage.\section{Overlapping Generations}The `selovlp.f90` module handles selection in age-structured populations. The key difference from the discrete generation model is that the genetic gain is calculated per year, rather than per generation.\subsection{Genetic Gain}The annual genetic gain is a function of the generation interval ($L$) and the genetic gain per generation ($\Delta G_{gen}$):\begin{equation}\Delta G_{year} = \frac{\Delta G_{gen}}{L}\end{equation}The generation interval is the average age of parents when their offspring are born, weighted by the number of offspring. The program calculates the generation interval for sires ($L_s$) and dams ($L_d$) separately and then averages them.\subsection{Gene Flow Method}The program uses a gene flow matrix method to model the flow of genes between different age classes over time. This allows for the calculation of genetic gain and inbreeding in complex population structures with multiple age classes for both sires and dams.\section{Inbreeding}The rate of inbreeding ($\Delta F$) is calculated using the approach of Bijma and Woolliams (2000), which is implemented in `selinbreeding.f90`. This method predicts the rate of inbreeding from BLUP selection based on long-term genetic contributions. The rate of inbreeding is given by:\begin{equation}\Delta F = \frac{1}{2} \left( \frac{1}{N_m} \sigma^2_{c,m} + \frac{1}{N_f} \sigma^2_{c,f} \right)\end{equation}where $N_m$ and $N_f$ are the number of male and female parents, and $\sigma^2_{c,m}$ and $\sigma^2_{c,f}$ are the variances of the long-term genetic contributions of males and females, respectively. These variances are a function of the selection process and the population structure.\end{document}